\subsection{The channels $(b,\partial_{\dot\alpha}c)$ and $(\partial_{\dot\alpha}c,b)$}\label{sec:superspace:ad}

We now compute the mixed-chirality family in full: the ordered channels
$(b,\partial_{\dot\alpha}c)$ and $(\partial_{\dot\alpha}c,b)$,
$\dot\alpha=\dot1,\dot2$, the block $4_{(b,\partial c),(\partial c,b)}$ of the
census~\eqref{eq:ss:29}. The computation repeats the seed-channel pattern of
Section~\ref{sec:superspace:seed} with three new features: the cubic vertices
have opposite chirality, the $D$-algebra numerator is rank one in the loop
momentum, and the ordered source admits two color allocations. Graph weights
are computed with unit-normalized insertions $\nabla_{+}\cW_{+}$ and
$\bcW_{\dot\alpha}$; the letter normalization of
Definition~\ref{def:ss:letters} is restored in step 7. The dotted transport
$\dot1\mapsto\dot2$ leaves $f^{ACD}f^{BCE}$ and the sign fixed, so we display
$\dot\alpha=\dot1$.

\subsubsection*{1. Graph inventory and the occurrence census}

Each ordering is carried by \emph{one ordered triangle plus one Schwinger
cut/contact graph}, paired by the cutting rule of
Section~\ref{sec:superspace:cutting}. The triangle
(Figure~\ref{fig:ss:ad-triangle}) has the ordered bilocal insertion $I$ and two
cubic background vertices of opposite chirality, the same mixed pair as the
seed chain,
\begin{equation}\label{eq:ss:ad-vertices}
V_{\widetilde\cW}:\ +\frac{ig}{2}\,f^{UCD}\,
\bcW^{D}_{\dot\gamma}\,\big(\bcD_{C}^{\dot\gamma}-\bcD_{U}^{\dot\gamma}\big),
\qquad
V_{\cW}:\ -\frac{ig}{2}\,f^{VC'E}\,
\cW^{E}_{\gamma}\,\big(\cD_{C'}^{\gamma}-\cD_{V}^{\gamma}\big).
\end{equation}
$V_{\cW}$ emits the $\nabla_{+}\cW_{+}$ insertion (the $b$ leg),
$V_{\widetilde\cW}$ emits the $\bcW_{\dot\alpha}$ insertion (the
$\partial_{\dot\alpha}c$ leg), and the three internal vector lines carry the
raw routing $r_{0}=\ell$, $r_{1}=\ell+p$, $r_{2}=\ell+p+q$, so
$D_{0}=\ell^{2}$, $D_{1}=(\ell+p)^{2}$, $D_{2}=(\ell+p+q)^{2}$ as in the seed
channel. The marked line is $e_{0}$ for $(b,\partial_{\dot\alpha}c)$ and
$e_{2}$ for $(\partial_{\dot\alpha}c,b)$.

\begin{figure}[t]
\centering
\begin{tikzpicture}[scale=1.2]
  \node[svertex, label=below left:{$I$}] (I) at (0,0) {};
  \node[svertex, label=above:{$V_{\widetilde\cW}$}] (Vt) at (2.9,1.7) {};
  \node[svertex, label=below:{$V_{\cW}$}] (Vw) at (2.9,-1.7) {};
  \draw[sV] (I) -- (Vt) node[midway, above left=-2pt, font=\small] {$e_{0}{:}\ \ell$};
  \draw[sV] (Vt) -- (Vw) node[midway, right, font=\small] {$e_{1}{:}\ \ell+p$};
  \draw[sV] (Vw) -- (I) node[midway, below left=-2pt, font=\small] {$e_{2}{:}\ \ell+p+q$};
  \draw[sV] (Vt) -- ++(1.9,0.95) node[midway, sinsert,
        label=above right:{$\partial_{\dot\alpha}c=i\bcW_{\dot\alpha}$}] {};
  \draw[sV] (Vw) -- ++(1.9,-0.95) node[midway, sinsert,
        label=below right:{$b=-\frac{i}{\sqrt2}\nabla_{+}\cW_{+}$}] {};
  \node[left=4pt of I, font=\small, align=right] {ordered bilocal\\ insertion $J_{AB}$};
\end{tikzpicture}
\caption{The ordered triangle of $(b,\partial_{\dot\alpha}c)$ and
$(\partial_{\dot\alpha}c,b)$. The bilocal insertion $I$ carries the two
letters; $V_{\cW}$ emits the $\nabla_{+}\cW_{+}$ insertion (the $b$ leg) and
$V_{\widetilde\cW}$ the $\bcW_{\dot\alpha}$ insertion (the
$\partial_{\dot\alpha}c$ leg); all internal lines are vector-superfield
propagators. The marked line is $e_{0}$ for $(b,\partial_{\dot\alpha}c)$ and
$e_{2}$ for $(\partial_{\dot\alpha}c,b)$; the Schwinger cut/contact partner
collapses the marked line with the same two vertices.}
\label{fig:ss:ad-triangle}
\end{figure}

The raw engine census per ordering is $36$ TGG $+$ $6$ TMM $+$ $42$ spectator
$=84$ connected $=42_{\rm split\ source}+42_{\rm spectator}$, where TGG (two
gauge cubic vertices) is one ordered triangle of
Figure~\ref{fig:ss:ad-triangle} with a distinguished marked edge, and TMM has
two matter vertices. The occurrence resolution: the $36$ TGG rows per
ordering, with their same-edge Schwinger cuts, carry the entire anomaly (steps
2--6). The $6$ TMM rows vanish exactly (wrong external field type for two
gauge letters). The matter-current Euler and explicit-contact pair shares the
word $K_{\dot1}=ir_{+\dot2}$, $K_{\dot2}=-ir_{+\dot1}$ with weights
$\pm2iK_{\dot\alpha}$ and cancels before regulation; the word is affine, so
its $\mul^{2}$ coefficient vanishes. The source-tagged bubbles have affine
numerators and zero standalone anomaly; their two-denominator rows belong to
the inverse-kernel brackets of step 4 and are not additional finite diagrams.
The tadpole is scaleless, $\int d^{d}\ell\,P(\ell)/\ell^{2}=0$ in DRED
\cite{SiegelDRED}. The gauge-fixing Euler sector cancels the longitudinal
source row (step 4) on $32$ component masks. In the ghost sector, the Euler
bubble and its full-$d$ cut are exact negatives and the ghost kernel is the
full-$d$ $\BoxE$, so no $\mul^{2}$ residue arises; the ghost triangle with
linear source is absent at one loop by the edge/vertex count, and the
Nielsen--Kallosh factor is absent from the selected graph set. The $42$
spectator routes per ordering are connected but one-particle reducible
(external-leg self-energies), and the amputated 1PI projector annihilates
them.

\begin{remark}[Status of the auxiliary sectors]\label{rem:ss:ad-census}
The local Fermi--Feynman proper slice is not derived from first principles; it
is fixed by the scheme decisions of Section~\ref{sec:gaugebv} (Gaussian
averaging; the Nielsen--Kallosh factor kept external). Every contributing
graph is accounted for by the triangle and its same-edge cut; all other
occurrences are exactly zero, cancel in pairs, or are absent.
\end{remark}

\subsubsection*{2. Wick contractions}

The two vertices of~\eqref{eq:ss:ad-vertices} are distinct (one chiral, one
antichiral), so the interaction expansion contributes once at second order:
\begin{equation}\label{eq:ss:ad-wick}
w_{\mathrm{Wick}}=\frac{1}{2!}\,(1+1)=1,
\end{equation}
the $(1+1)$ counting the two equivalent Wick pairings, and the labeled source
Hessian contributes $1!\,1!=1$. There is \emph{no reflection multiplicity}:
$(b,\partial_{\dot\alpha}c)$ (marked edge $e_{0}$) and
$(\partial_{\dot\alpha}c,b)$ (marked edge $e_{2}$) are distinct ordered
descendants, stored separately; the compact output word is already oriented.
The three internal lines contract against the vertices and the source with the
vector propagators of the Fermi--Feynman slice ($\kappa_{AB}=\delta_{AB}$,
Section~\ref{sec:gaugebv}):
\begin{equation}\label{eq:ss:ad-props}
\langle\cV^{A}\cV^{U}\rangle_{e_{0}}
=\frac{\h\,\delta^{AU}\delta^{4}(\theta_{12})}{\ell^{2}},
\qquad
\langle\cV^{C}\cV^{C'}\rangle_{e_{1}}
=\frac{\h\,\delta^{CC'}\delta^{4}(\theta_{12})}{(\ell+p)^{2}},
\qquad
\langle\cV^{V}\cV^{B}\rangle_{e_{2}}
=\frac{\h\,\delta^{VB}\delta^{4}(\theta_{12})}{(\ell+p+q)^{2}} .
\end{equation}
\begin{remark}[Edge-weight sign convention]\label{rem:ss:ad-edgesign}
The edge weights displayed in~\eqref{eq:ss:ad-props} are recorded in the
graph-IR convention of the underlying computation, in which the overall
minus sign of the Feynman propagator $\langle uu\rangle=-\h\kappa^{AB}/p^{2}$
of Section~\ref{sec:superspace:rules} is folded into the vertex signs of the
closed word; the displayed $+\h\kappa^{AB}/p^{2}$ is the edge weight per
line, not an independent propagator sign. All downstream factors of this
subsection use the same convention consistently, and the final coefficients
are machine-verified against the propagator-level replay.
\end{remark}
Contraction with the vertex structure constants and the source yields the
common color tensor in its two allocations, $f^{ACD}f^{BCE}$ (direct) and
$f^{BCD}f^{ACE}$ (crossed), exchanged by $A\leftrightarrow B$ and independent
for a general gauge algebra. The surviving direct ordered source Hessians on
the fixed odd component directions are $i/4$ for $(b,\partial_{\dot\alpha}c)$
and $1/4$ for $(\partial_{\dot\alpha}c,b)$, with $1!\,1!=1$; each external
$\partial_{\dot\alpha}c$ leg carries the component preimage
$-\frac{\sqrt2}{8}\eta_{\dot\alpha}$. Reversing the functional order of the
two odd directions changes the sign of the crossed Hessian, which is why both
color allocations must be kept.

\subsubsection*{3. Amplitude assembly}

The two ordered cubic vertices contribute
$\big(+\frac{ig}{2}\big)\big(-\frac{ig}{2}\big)=\frac{g^{2}}{4}$
as in~\eqref{eq:ss:vertices}, and the closed loop is saturated by
$\delta^{4}(\theta_{12})\cD^{2}\bcD^{2}\delta^{4}(\theta_{12})
=16\,\delta^{4}(\theta_{12})$. The expansion runs over the four chirality
sectors $(++),(+-),(-+),(--)$ of the two vertices, each of weight $1$; the
four contractions coincide polynomially, so the full action contribution is
$4_{\rm chirality}\times(\text{one sector})$. For each occurrence $(\omega,e)$
the frozen source identity splits the marked $D$-algebra word as
\begin{equation}\label{eq:ss:ad-split}
\cD_{-}\cD_{+}\bcD^{2}\cD_{+}
=-8\,\bar r_{e}^{\,2}\,\cD_{+}-\frac12\,\cD_{+}\bcD^{2}\cD^{2},
\end{equation}
so that, writing $\mathscr W^{K}_{\omega,e}$ and $\mathscr W^{L}_{\omega,e}$
for the kinetic and longitudinal spin words,
\begin{equation}\label{eq:ss:ad-kl}
K_{\omega,e}=-8\,\bar r_{e}^{\,2}\,\mathscr W^{K}_{\omega,e},
\qquad
L^{\mathrm{source}}_{\omega,e}=-\tfrac12\,\mathscr W^{L}_{\omega,e},
\qquad
L^{\mathrm{gf}}_{\omega,e}=+\tfrac12\,\mathscr W^{L}_{\omega,e},
\qquad
L^{\mathrm{source}}_{\omega,e}+L^{\mathrm{gf}}_{\omega,e}=0 .
\end{equation}
With $R_{\omega,e}$ the reduced numerator and
$C^{\mathrm{SD}}_{\omega,e}=-R_{\omega,e}/\prod_{j\ne e}D_{j}$ the same-edge
Schwinger contact, the full-$d$ and four-dimensional identities are
\begin{equation}\label{eq:ss:ad-sd}
K^{(d)}_{\omega,e}+C^{\mathrm{SD}}_{\omega,e}=0,
\qquad
K^{(4)}_{\omega,e}+C^{\mathrm{SD}}_{\omega,e}
=\frac{\mul^{2}\,R_{\omega,e}}{D_{0}D_{1}D_{2}},
\end{equation}
the edgewise parent-minus-cut of~\eqref{eq:ss:cuttingfailure}.

\begin{remark}[The factored weight chain is recorded, not displayed]%
\label{rem:ss:ad-weightchain}
Unlike the seed channel, no single displayed chain
$w_{\mathrm{Wick}}\frac{g^{2}}{4}w_{D}$ exists for this family; the total is
fixed by the recorded conversions: a per-chirality conditional factor $-8$, the
absolute conversion
$4_{\rm chirality}\times\big(\tfrac12\big)_{\rm rank/trace}\times
(-8)_{\rm legacy}=-16$,
and the isolating color value $-2$ (the $SU(2)$ sample $(A,B;D,E)=(0,1;1,0)$
gives $(f^{ACD}f^{BCE},f^{BCD}f^{ACE})=(-2,0)$). The rank/trace factor
$\tfrac12$ is inserted exactly once; $m_{\mathrm{contact}}=1$ of step 6
supplies no second factor.
\end{remark}

\subsubsection*{4. The $D$-algebra numerator chain}

With the conditional Fermi--Feynman projector
$\mathcal P_{0}=(\bcD^{2}\cD^{2}+\cD^{2}\bcD^{2})/(16\BoxE)$ and
$\cD_{+}\cD^{2}\bcD^{2}=0$,
\begin{equation}\label{eq:ss:ad-projector}
8\,\cD_{+}\BoxE\,\mathcal P_{0}
=\tfrac12\,\cD_{+}\bcD^{2}\cD^{2}+\tfrac12\,\cD_{+}\cD^{2}\bcD^{2}
=\tfrac12\,\cD_{+}\bcD^{2}\cD^{2},
\end{equation}
so the longitudinal term of~\eqref{eq:ss:ad-split} is pure gauge and cancels
the gauge-fixing Euler row,
$-\tfrac12\cD_{+}\bcD^{2}\cD^{2}+\tfrac12\cD_{+}\bcD^{2}\cD^{2}=0$, verified
sparsely on $32+32$ component masks. What survives is the first term
of~\eqref{eq:ss:ad-split}: per marked edge the parent-minus-cut
of~\eqref{eq:ss:ad-sd} leaves exactly the evanescent
remainder~\eqref{eq:ss:cuttingfailure}. Summing the $36$ TGG routes per
ordering gives the two rank-one route sums
\begin{equation}\label{eq:ss:ad-r12}
R_{1}=\frac{i}{16}\,\mul^{2}\,\Big(\ell+p+\frac12q\Big)_{+\dot1},
\qquad
R_{2}=-\frac{i}{16}\,\mul^{2}\,\ell_{+\dot1},
\end{equation}
with $X_{+\dot1}$ the plus-frame dotted spinor component, and the evanescent
numerators of the two orderings carry the color mix
\begin{equation}\label{eq:ss:ad-nev}
N^{\mathrm{ev}}_{(b,\partial c)}
=f^{ACD}f^{BCE}\,R_{1}+f^{BCD}f^{ACE}\,R_{2},
\qquad
N^{\mathrm{ev}}_{(\partial c,b)}
=f^{ACD}f^{BCE}\,R_{2}+f^{BCD}f^{ACE}\,R_{1},
\end{equation}
the marked edge being $e_{0}$ for the direct $(b,\partial c)$ and crossed
$(\partial c,b)$ rows and $e_{2}$ for the other two. Occurrence by occurrence,
the aggregated two-denominator contact rows and the same-edge remainders obey
\begin{align}
(b,\partial c),\ e_{0}:&\quad
7iq_{+\dot1}+\Big[i\Big(\ell+p+\tfrac12q\Big)_{+\dot1}-7iq_{+\dot1}\Big]
=i\Big(\ell+p+\tfrac12q\Big)_{+\dot1},\notag\\
(b,\partial c),\ e_{1}:&\quad 0+0=0,
\qquad\qquad
(b,\partial c),\ e_{2}:\quad 2iq_{+\dot1}-2iq_{+\dot1}=0,\notag\\
(\partial c,b),\ e_{0}:&\quad
-\tfrac{i}{2}q_{+\dot1}+\tfrac{i}{2}q_{+\dot1}=0,
\qquad\quad
(\partial c,b),\ e_{1}:\quad
i(p+q)_{+\dot1}-i(p+q)_{+\dot1}=0,\notag\\
(\partial c,b),\ e_{2}:&\quad
i(p+q)_{+\dot1}+\big[-i\ell_{+\dot1}-i(p+q)_{+\dot1}\big]
=-i\ell_{+\dot1},
\label{eq:ss:ad-sameedge}
\end{align}
and the crossed-color replay exchanges the two orderings edge by edge. The
two-denominator members of~\eqref{eq:ss:ad-sameedge} belong to the
inverse-kernel brackets and must not be added again as finite diagrams once
$N^{\mathrm{ev}}$ is formed.

\subsubsection*{5. Tensor integrals and the ordered simplex moments}

The two DRED masters are the scalar master~\eqref{eq:ss:masterintegral} and
its rank-one partner,
\begin{equation}\label{eq:ss:ad-masters}
\int\!\frac{d^{d}\ell}{(2\pi)^{d}}\,\frac{\mul^{2}}{D_{0}D_{1}D_{2}}
=\frac{1}{32\pi^{2}},
\qquad
\int\!\frac{d^{d}\ell}{(2\pi)^{d}}\,
\frac{\mul^{2}\,\ell_{+\dot1}}{D_{0}D_{1}D_{2}}
=-\frac{(2p+q)_{+\dot1}}{96\pi^{2}} .
\end{equation}
Combining denominators,
\begin{equation}\label{eq:ss:ad-feynman}
\frac{1}{D_{0}D_{1}D_{2}}
=2\int_{0}^{1}\!dy\int_{0}^{1-y}\!dz\,\frac{1}{(L^{2}+\Delta)^{3}},
\qquad
L=\ell+(y+z)p+zq,
\end{equation}
with exact simplex moments, writing $x=1-y-z$,
\begin{equation}\label{eq:ss:ad-simplex}
2\!\int_{\Delta_{2}}\!1=1,\qquad
2\!\int_{\Delta_{2}}\!x=\frac13,\qquad
2\!\int_{\Delta_{2}}\!(y+z)=\frac23,\qquad
2\!\int_{\Delta_{2}}\!z=\frac13 .
\end{equation}
For $(b,\partial_{\dot\alpha}c)$ the shifted numerator is
$\ell+p+\frac12q=L+xp+\big(\frac12-z\big)q$; the term odd in $L$ integrates to
zero, and $2\int_{\Delta_{2}}(\tfrac12-z)=2(\tfrac14-\tfrac16)=\tfrac16$, so
\begin{equation}\label{eq:ss:ad-adint}
\int_{\ell}\frac{N^{\mathrm{ev}}_{(b,\partial c)}}{D_{0}D_{1}D_{2}}\Big|_{R_{1}}
=\frac{i}{16}\cdot\frac{1}{32\pi^{2}}
\left(\frac13\,p+\frac16\,q\right)_{+\dot1}
=\frac{i}{512\pi^{2}}\left(\frac13\,p+\frac16\,q\right)_{+\dot1} .
\end{equation}
For $(\partial_{\dot\alpha}c,b)$, $-\ell=-L+(y+z)p+zq$ and the
moments~\eqref{eq:ss:ad-simplex} give
\begin{equation}\label{eq:ss:ad-daint}
\int_{\ell}\frac{N^{\mathrm{ev}}_{(\partial c,b)}}{D_{0}D_{1}D_{2}}\Big|_{R_{2}}
=-\frac{i}{16}\left(-\frac{1}{32\pi^{2}}\right)
\left(\frac23\,p+\frac13\,q\right)_{+\dot1}
=\frac{i}{512\pi^{2}}\left(\frac23\,p+\frac13\,q\right)_{+\dot1} ,
\end{equation}
in agreement with the direct evaluation of~\eqref{eq:ss:ad-masters},
$-\frac{i}{16}\big(-\frac{(2p+q)_{+\dot1}}{96\pi^{2}}\big)
=\frac{i}{512\pi^{2}}(\frac23p+\frac13q)_{+\dot1}$. In the authority routing
the raw external momenta are the ordered action ports, $p=k_{L}$ (left output
leg) and $q=k_{R}$ (right output leg). The jet weights of the final bilinear
are the ordered-simplex moments of the tower
parametrization~\eqref{eq:ss:simplexderivation} at $(m,n)=(1,0)$:
\begin{align}
2\int_{0}^{1}\!db\int_{0}^{b}\!da\;a
&=2\int_{0}^{1}\frac{b^{2}}{2}\,db=\frac13
\quad(\text{jet on the first letter}),\notag\\
2\int_{0}^{1}\!db\int_{0}^{b}\!da\;b
&=2\int_{0}^{1}b^{2}\,db=\frac23
\quad(\text{jet on the second letter}),
\label{eq:ss:ad-simplex10}
\end{align}
precisely $K_{1,0;1,0}=\frac{2}{3\cdot2}$ and $K_{1,0;0,0}=\frac{2}{3\cdot1}$
of~\eqref{eq:ss:kernel}. Assembling the four chirality sectors, the two color
allocations and the contact rows, the target-blind momentum shapes are
\begin{align}
\Gamma_{(b,\partial c)}
&=-i\,\frac{\h g^{2}}{16\pi^{2}}\left[
f^{ACD}f^{BCE}\left(\frac13k_{L}+\frac16k_{R}\right)
+f^{BCD}f^{ACE}\left(\frac23k_{L}+\frac13k_{R}\right)\right]_{+\dot1},\notag\\
\Gamma_{(\partial c,b)}
&=-i\,\frac{\h g^{2}}{16\pi^{2}}\left[
f^{ACD}f^{BCE}\left(\frac23k_{L}+\frac13k_{R}\right)
+f^{BCD}f^{ACE}\left(\frac13k_{L}+\frac16k_{R}\right)\right]_{+\dot1}.
\label{eq:ss:ad-shapes}
\end{align}

\begin{remark}[The jet readout is the assembled output, not a Fourier map]%
\label{rem:ss:ad-fourier}
A naive identification of momentum shape with jet weights fails: the
direct-color shape in~\eqref{eq:ss:ad-shapes} is $(\tfrac13,\tfrac16)$, not
$(\tfrac13,\tfrac23)$. The canonical bilinear coefficients of step 7 are the
recorded output of the full four-chirality, two-color, contact-inclusive
assembly, not of a one-line Fourier identification;
\eqref{eq:ss:ad-simplex10} displays the intrinsic jet weights of that
assembled result.
\end{remark}

\subsubsection*{6. Pole cancellation}

The triangle parent and the aggregate Schwinger contact row carry the UV-pole
pair
\begin{align}
\Gamma_{T}&=+\frac{\h g^{2}}{32\pi^{2}\epsilon}\times(\text{edge-tagged
tensor with }\widehat\delta^{mn}),\notag\\
\Gamma_{C}&=-\frac{\h g^{2}}{32\pi^{2}\epsilon}\times(\text{the identical
tensor with }\delta_{4}^{mn}),
\label{eq:ss:ad-polepair}
\end{align}
so the four-dimensional metric parts cancel exactly and the remainder is the
pure evanescent numerator $\widehat\delta^{mn}-\delta_{4}^{mn}
=-\breve\delta^{mn}$, whose sigma-trace supplies the $-2\epsilon$
of~\eqref{eq:ss:sigmatrace}: finiteness holds precisely as in
Remark~\ref{rem:ss:polebookkeeping}. At the occurrence level, the local
contact kernel $K_{AB}=-\kappa_{AB}r_{e}^{2}$ with
$\kappa_{AB}=\tfrac12\delta_{AB}$ obeys $r_{e}^{2}P^{A}=-2P^{A}K_{A}$, and
\begin{align}
c_{\mathrm{parent}}&=(-2)^{3}=-8,
\qquad
c_{\mathrm{contact}}
=(-1)_{\mathrm{SD}}\,(-2)_{r^{2}P=-2PK}\,(-2)^{2}_{\mathrm{uncut}}=+8,\notag\\
m_{\mathrm{contact}}
&=\Big(\tfrac{1}{2!}\,2_{\mathrm{action}}\Big)
\Big(\tfrac{1}{2!}\,2_{\mathrm{cut}}\Big)
\Big(\tfrac{1}{2!}\,2_{\mathrm{Wick}}\Big)=1 .
\label{eq:ss:ad-signs}
\end{align}
Consequently, occurrence by occurrence and in each chirality sector,
\begin{equation}\label{eq:ss:ad-occcancel}
N^{\mathrm{parent}}_{d,\omega}+K^{\mathrm{contact}}_{\mathrm{raw},\omega}=0,
\qquad
N^{\mathrm{parent}}_{4,\omega}+K^{\mathrm{contact}}_{\mathrm{raw},\omega}
=\frac{R_{\omega}\,\mul^{2}}{D_{0}D_{1}D_{2}} .
\end{equation}
The finite-$w$ link sector cancels before loop integration: with $x=iwr$,
$y=iwr'$ and $(x-y)\int_{0}^{1}ds\,e^{sx+(1-s)y}=e^{x}-e^{y}$, the one-link
term and the two endpoint terms cancel,
$C_{\omega,e}(e^{x}-e^{y})-C_{\omega,e}(e^{x}-e^{y})=0$; the sector is
longitudinal in arbitrary $d$, so its $\mul^{2}$ coefficient is exactly zero.

\subsubsection*{7. The final coefficients and the tower-kernel instance}

With unit-normalized insertions the assembled channel results are
\begin{align}
\nabla_{-}\big((\nabla_{+}\cW_{+})^{A}\,\bcW_{\dot\alpha}^{B}\big)
&=\frac{\h g^{2}}{16\pi^{2}}\,f^{ACD}f^{BCE}\,
\Big\{\tfrac13\big\langle\partial_{\dot\alpha}\bcW,\bcW\big\rangle
+\tfrac23\big\langle\bcW,\partial_{\dot\alpha}\bcW\big\rangle\Big\}^{DE},
\label{eq:ss:ad-unit1}\\
\nabla_{-}\big(\bcW_{\dot\alpha}^{A}\,(\nabla_{+}\cW_{+})^{B}\big)
&=\frac{\h g^{2}}{16\pi^{2}}\,f^{ACD}f^{BCE}\,
\Big\{\tfrac23\big\langle\partial_{\dot\alpha}\bcW,\bcW\big\rangle
+\tfrac13\big\langle\bcW,\partial_{\dot\alpha}\bcW\big\rangle\Big\}^{DE}.
\label{eq:ss:ad-unit2}
\end{align}
Restoring the letters of Definition~\ref{def:ss:letters} is exact arithmetic:
the insertions and the output letters give
\begin{align}
b^{A}\,\partial_{\dot\alpha}c^{B}
&=\Big(-\frac{i}{\sqrt2}\Big)(i)\,
(\nabla_{+}\cW_{+})^{A}\,\bcW_{\dot\alpha}^{B}
=\frac{1}{\sqrt2}\,(\nabla_{+}\cW_{+})^{A}\,\bcW_{\dot\alpha}^{B},
\notag\\
\big\langle\partial_{\dot\alpha}\bcW,\bcW\big\rangle
&=(-i)^{2}\,\big\langle\partial_{\dot\alpha}(\partial c),\partial c\big\rangle ,
\label{eq:ss:ad-restore}
\end{align}
and $\frac{1}{\sqrt2}\times(-1)\times\frac{\h g^{2}}{16\pi^{2}}
=-\frac{\h g^{2}}{16\sqrt2\,\pi^{2}}$. Substituting
into~\eqref{eq:ss:ad-unit1}--\eqref{eq:ss:ad-unit2},
\begin{equation}\label{eq:ss:ad-bdc}
\boxed{\begin{aligned}
\nabla_{-}\big(b^{A}\partial_{\dot\alpha}c^{B}\big)
&=-\frac{\h g^{2}}{16\sqrt2\,\pi^{2}}\,f^{ACD}f^{BCE}\\
&\quad\times
\Big\{\tfrac13\big\langle\partial_{\dot\alpha}(\partial c),\partial c\big\rangle
+\tfrac23\big\langle\partial c,\partial_{\dot\alpha}(\partial c)\big\rangle
\Big\}^{DE}
\end{aligned}}
\end{equation}
\begin{equation}\label{eq:ss:ad-dcb}
\boxed{\begin{aligned}
\nabla_{-}\big(\partial_{\dot\alpha}c^{A}b^{B}\big)
&=-\frac{\h g^{2}}{16\sqrt2\,\pi^{2}}\,f^{ACD}f^{BCE}\\
&\quad\times
\Big\{\tfrac23\big\langle\partial_{\dot\alpha}(\partial c),\partial c\big\rangle
+\tfrac13\big\langle\partial c,\partial_{\dot\alpha}(\partial c)\big\rangle
\Big\}^{DE}
\end{aligned}}
\end{equation}
--- exactly the rows~\eqref{eq:ss:ch-bdc} and~\eqref{eq:ss:ch-dcb} of the
channel table. The weights are the $(m,n)=(1,0)$ instance of the tower
kernel~\eqref{eq:ss:kernel}: by~\eqref{eq:ss:ad-simplex10},
$K_{1,0;1,0}=\frac{2}{3\cdot2}=\frac13$ (jet on the first letter) and
$K_{1,0;0,0}=\frac{2}{3\cdot1}=\frac23$ (jet on the second), and the two
orderings are exchanged by the Leibniz distribution
$(1,1)-(\tfrac13,\tfrac23)=(\tfrac23,\tfrac13)$.

\begin{remark}[No $1:2$ ratio between the two orderings]\label{rem:ss:ad-ratio}
An earlier conditional proposal assigned the two orderings the ratio
$\mathcal N_{(b,\partial c)}/\mathcal N_{(\partial c,b)}=1/2$ through a
rejected momentum map; it was explicitly \emph{not} promoted to a coefficient
of this paper. The final results~\eqref{eq:ss:ad-bdc}--\eqref{eq:ss:ad-dcb}
share the single prefactor
$-\frac{\h g^{2}}{16\sqrt2\,\pi^{2}}f^{ACD}f^{BCE}$ and differ only in the
jet weights $(\tfrac13,\tfrac23)\leftrightarrow(\tfrac23,\tfrac13)$.
\end{remark}

\begin{remark}[Machine cross-checks for this family]\label{rem:ss:ad-checks}
The full chain of this subsection was replayed externally: $94/94$ checks on
the ordered-port crossed-Hessian assembly, $112/112$ on each of the two
polarized $D$-algebra replays, and $255/255$ on the raw contact and finite-$w$
link replay, with zero failures.
\end{remark}

\subsubsection*{8. Comparison with the holomorphic twist}

In the holomorphic twist of~\cite{BK} this family has no zero-derivative
display; it exists only through the derivative tower applied to the seed pair
$Q_{1}(b^{A}c^{B})
=\kappa_{H}^{2}f_{ACD}f_{BCE}\,\partial_{\dot\alpha}c^{D}
\partial^{\dot\alpha}c^{E}$ (their eq.~(3.63)). Their Appendix-B tower
(their eq.~(B.6)),
\begin{align}
Q_{1}\big(b^{A}\partial_{1}^{m}\partial_{2}^{n}c^{B}\big)
&=\kappa_{H}^{2}f_{ACD}f_{BCE}\,\frac{1}{m+n+2}
\sum_{k=0}^{m}\sum_{\ell=0}^{n}\frac{1}{k+\ell+1}\binom{m}{k}\binom{n}{\ell}
\notag\\
&\quad\times
\big[\partial_{1}^{k}\partial_{2}^{\ell}\partial_{\dot\alpha}c^{D}\,
\partial_{1}^{m-k}\partial_{2}^{\,n-\ell}\partial^{\dot\alpha}c^{E}\big],
\label{eq:ss:ad-bktower}
\end{align}
evaluated at $(m,n)=(1,0)$, gives the printed weights
\begin{align}
k=0:&\quad \frac13\ \text{on}\ \partial_{\dot\alpha}c^{D}\,
\partial_{1}\partial^{\dot\alpha}c^{E}\ (\text{jet on the second letter}),
\notag\\
k=1:&\quad \frac16\ \text{on}\ \partial_{1}\partial_{\dot\alpha}c^{D}\,
\partial^{\dot\alpha}c^{E}\ (\text{jet on the first letter}),
\label{eq:ss:ad-bkweights}
\end{align}
exactly one half of our coefficients on both words. After the factor-two
adjudication of Section~\ref{sec:ht} --- the printed Appendix-B branch
of~\cite{BK} is uniformly smaller by the factor $2=\Gamma(3)$
of~\eqref{eq:ss:feynman} --- the corrected weights $(\tfrac13,\tfrac23)$
coincide with~\eqref{eq:ss:ad-bdc}. The reversed ordering
$(\partial_{\dot\alpha}c,b)$ is not printed in~\cite{BK} at all; it follows
from their recursion (their eq.~(B.1)),
$Q_{1}(\partial_{\dot\alpha}fg)
=\partial_{\dot\alpha}Q_{1}(fg)-Q_{1}(f\,\partial_{\dot\alpha}g)$, applied to
$(f,g)=(c,b)$, the unprinted orderings being identified with the printed ones
by supercommutativity of the letters (no sign, since $|c||b|=0$). On our side
the same law is the Leibniz distribution on the formal ghost slot: the seed
distributes one unit of weight on each output word, $(1,1)$, and
$(1,1)-(\tfrac13,\tfrac23)=(\tfrac23,\tfrac13)$, exactly the swap
between~\eqref{eq:ss:ad-bdc} and~\eqref{eq:ss:ad-dcb}. Both dotted rows agree
word by word with the twist computation (Section~\ref{sec:ht}).

\begin{remark}[Status of the ghost-slot recursion]\label{rem:ss:ad-Uslot}
The recursion anchor involves the formal odd slot of the twist dictionary (the
ghost placeholder, of parity $1$), whose identification with project
quantities is conditional; the physical
rows~\eqref{eq:ss:ad-bdc}--\eqref{eq:ss:ad-dcb} are derived independently in
steps 1--7 and are unaffected. The comparison of this subsection is a
cross-check, not a derivation input, and not a completeness proof of the graph
census (Section~\ref{sec:ht}).
\end{remark}
